% ---------------------------------------------------------------------
% Parte 4 · Ejemplo sencillo · expone César · 2:00 (lo edita solo César)
% El código del ejemplo va en ejemplo/ (con su README de cómo ejecutarlo).
% En la segunda diapositiva, reemplacen el pendiente por las 10-12 líneas
% clave del ejemplo, por ejemplo:
%   \lstinputlisting[language=Java, firstline=5, lastline=16]{ejemplo/Servidor.java}
% Tengan un video o capturas de respaldo por si la demo falla en clase.
% ---------------------------------------------------------------------
\begin{frame}{Qué vamos a ver}
  \vspace{0.3cm}
  \begin{center}
    \begin{tikzpicture}[x=1cm, y=1cm]
      \foreach \n/\x in {1/0, 2/4.6, 3/9.2} {
        \fill[bcmWS] (\x,0) circle (0.42);
        \node[text=white, font=\large\bfseries] at (\x,0) {\n};
      }
      \draw[bcmGris, line width=0.8pt, -stealth] (0.6,0) -- (4.0,0);
      \draw[bcmGris, line width=0.8pt, -stealth] (5.2,0) -- (8.6,0);
      \foreach \x in {0, 4.6, 9.2}
        \node[anchor=north, text width=3.8cm, align=center, font=\small] at (\x,-0.65)
          {\pend{César}{paso que verán en pantalla}};
    \end{tikzpicture}
  \end{center}
  \vfill
  {\small\color{bcmGris}\pend{César}{con qué está hecho el ejemplo y dónde está en el repositorio}\par}
\end{frame}

\begin{frame}{El código que importa}
  \begin{columns}[T]
    \begin{column}{0.58\textwidth}
      \pendiente{César}{las 10 a 12 líneas clave del ejemplo (ver el comentario al inicio de este
        archivo)}
    \end{column}
    \begin{column}{0.38\textwidth}
      \begin{tarjeta}{bcmWS}{Fíjense en}
        \small
        \begin{itemize}
          \item \pend{César}{línea o idea 1}
          \item \pend{César}{línea o idea 2}
        \end{itemize}
      \end{tarjeta}
    \end{column}
  \end{columns}
\end{frame}
